% =============================================================================
%  Spoločné nastavenie všetkých troch dokumentov (preklad: latexmk -xelatex)
%  Záhrada stratených ciest – izometrické bludisko v knihe, Mykyta Olym, PG 2026/2027
% =============================================================================

% ---------------------------------------------------------------- údaje projektu
\newcommand{\nazovprojektu}{Záhrada stratených ciest}
\newcommand{\podnazovprojektu}{izometrické bludisko v knihe}
\newcommand{\predmet}{Počítačová grafika}
\newcommand{\skolskyrok}{2026/2027}
\newcommand{\ucitel}{Course Instructor}
\newcommand{\autor}{Mykyta Olym}
\newcommand{\verzia}{2.0}

% ---------------------------------------------------------------- jazyk a písmo
\usepackage{fontspec}
\usepackage{polyglossia}
\setdefaultlanguage{slovak}
\setmainfont{TeX Gyre Pagella}
\setsansfont{TeX Gyre Heros}
% neproporcionálne písmo zo šablóny KPI (priečinok fonts/)
\setmonofont{SourceCodePro}[
  Path = fonts/, Extension = .otf, Scale = 0.85,
  UprightFont = *-Regular, BoldFont = *-Bold,
  ItalicFont = *-Italic, BoldItalicFont = *-BoldItalic]

% ---------------------------------------------------------------- strana
\usepackage[a4paper, top=25mm, bottom=30mm, left=28mm, right=25mm]{geometry}
\usepackage{fancyhdr}
\pagestyle{fancy}
\fancyhf{}
\fancyhead[L]{\small\nazovprojektu}
\fancyhead[R]{\small\druhprirucky}
\fancyfoot[R]{\thepage}
\renewcommand{\headrulewidth}{0.4pt}
\setlength{\headheight}{14pt}
\setlength{\parindent}{0pt}
\setlength{\parskip}{0.55em}
\setlength{\emergencystretch}{3em}   % dlhé názvy radšej zalomiť než vytlačiť za okraj
% názvy súborov sa smú zalomiť za lomkou, bodkou alebo podčiarknutím
\usepackage[hyphens]{url}
\newcommand{\subor}[1]{\nolinkurl{#1}}

% ---------------------------------------------------------------- obsah, obrázky, tabuľky
\usepackage{graphicx}
\graphicspath{{obrazky/}}
\usepackage{booktabs, tabularx, array, longtable}
\usepackage{float}
\usepackage[font=small, labelfont=bf]{caption}
\usepackage{enumitem}
\setlist{itemsep=0.15em, topsep=0.3em}
\usepackage{xcolor}
\usepackage{csquotes}
\usepackage[style=iso-numeric, backend=biber, sorting=none]{biblatex}
\addbibresource{literatura.bib}
\usepackage[hidelinks]{hyperref}

% ---------------------------------------------------------------- titulná strana podľa pokynov
\newcommand{\titulnastrana}{%
  \begin{titlepage}
    \centering
    {\large Katedra počítačov a informatiky\\
      Fakulta elektrotechniky a informatiky\\
      Technická univerzita v Košiciach\par}
    \vspace*{\fill}
    {\Huge\bfseries \nazovprojektu\par}
    \vspace{0.4em}
    {\large \podnazovprojektu\par}
    \vspace{1.2em}
    {\Large \predmet\par}
    \vspace{3em}
    {\LARGE \druhprirucky\par}
    \vspace*{\fill}
    \begin{minipage}[t]{0.62\textwidth}
      \raggedright
      \begin{tabular}{@{}ll@{}}
        Šk. rok:       & \skolskyrok \\
        Meno učiteľa:  & \ucitel
      \end{tabular}
    \end{minipage}%
    \begin{minipage}[t]{0.38\textwidth}
      \raggedleft
      \begin{tabular}{@{}ll@{}}
        Autori: & \autor
      \end{tabular}
    \end{minipage}
  \end{titlepage}
  \setcounter{page}{1}%
}

% Obsah je podľa pokynov kapitola 1, preto skutočné kapitoly začínajú číslom 2.
\newcommand{\obsah}{%
  \renewcommand{\contentsname}{1\quad Obsah}%
  \tableofcontents
  \setcounter{section}{1}%
  \clearpage
}
